\documentclass{beamer}




\usepackage[T1,T2A]{fontenc}
\usepackage[utf8]{inputenc}

\usepackage{graphicx}
\usepackage{blindtext}
\usepackage{ mathrsfs }
\usepackage[russian,english]{babel}

\author{Деркач Максим Юрьевич}
\title{Криптографические протоколы}
\subtitle{Лекция 5 \\ Протоколы "сильной" аутентификации }


%\usetheme{lucid}
\begin{document}
	\frame {
		\titlepage
	}
	\frame {
		\frametitle{Ссылки}
		
	}
	
	\frame{
		\frametitle{Классификация протоколов "сильной" аутентификации}
		\begin{itemize}
				\item[*] на основе симметричных алгоритмов шифрования;
				\item[*] на основе асимметричных алгоритмов шифрования;
				\item[*] на основе ЭЦП.	
		\end{itemize}
	}

	\frame{
		\frametitle{Протоколы "сильной" аутентификации}
		\framesubtitle{На основе симметричных алгоритмов шифрования}
		
		\textbf{ISO/ IEC 9798 - 2}
		\begin{enumerate}
			\item $A -> B: text_2 || E_{K_{AB}}(T_A/N_A || ID_B || text_1)$
			\item $B -> A: text_4 || E_{K_{AB}}(T_B/N_B || ID_A || text_3)$
		\end{enumerate}

		\bigskip

		\begin{enumerate}
			\item $B -> A: R_B || text_1$
			\item $A -> B: text_3 || E_{K_{AB}}(R_B || ID_B || text_2)$
		\end{enumerate}
	
		\bigskip
		
		\begin{enumerate}
			\item $B -> A: R_B || text_1$
			\item $A -> B: text_3 || E_{K_{AB}}(R_A || R_B || ID_B || text_2)$
			\item $B -> A: text_5 || E_{K_{AB}}(R_B || R_A || text_4)$
		\end{enumerate}
		\textbf{Замечание} 
		Длина $ID_B || text_2$ не должна равняться длине $text_4$
		
		
	}
	
	\frame{
		\frametitle{Протоколы "сильной" аутентификации}
		\framesubtitle{На основе симметричных алгоритмов шифрования}
		
		\textbf{Протокол Ву-Лама (Woo-Lam)}
		\begin{enumerate}
			\item $A -> B: ID_A$
			\item $B -> A: R_B$
			\item $A -> B: E_{K_{AT}}(R_B)$
			\item $B -> T: E_{K_{BT}}(ID_A || E_{K_{AT}}(R_B))$
			\item $T -> B: E_{K_{BT}}(R_B)$
		\end{enumerate}
	}

	\frame{
		\frametitle{Протоколы "сильной" аутентификации}
		\framesubtitle{На основе симметричных алгоритмов шифрования}
		
		\textbf{Атака параллельного сеанса}
		\begin{enumerate}
			\item[1]  $I(A) -> B: ID_A$
			\item[1'] $I -> B: ID_I$
			
			\item[2]  $B -> I(A): R_B$
			\item[2'] $B -> I: R_B^*$
			
			
			\item[3]  $I(A) -> B: E_{K_{IT}}(R_B)$
			\item[3'] $I -> B: E_{K_{IT}}(R_B)$
			
			
			\item[4]  $B -> T: E_{K_{BT}}(ID_A || E_{K_{IT}}(R_B))$
			\item[4'] $B -> T: E_{K_{BT}}(ID_I || E_{K_{IT}}(R_B))$
			
			
			\item[5]  $T -> B: E_{K_{BT}}("MYCOP")$
			\item[5'] $T -> B: E_{K_{BT}}(R_B)$
		\end{enumerate}	
	}
	
		\frame{
		\frametitle{Протоколы "сильной" аутентификации}
		\framesubtitle{На основе симметричных алгоритмов шифрования}
		
		\textbf{Протокол Отвея-Риса}
		\begin{enumerate}
			\item $A -> B: ID_A || ID_B || K_{AT}(ID_A || ID_B || N_A || R_A) || N_A$
			\item $B -> T: ID_A || ID_B || K_{AT}(ID_A || ID_B || N_A || R_A) ||$\\ 
			 \quad \quad \quad  \quad \ $K_{BT}(ID_A || ID_B || N_A || R_B) || N_A$
			\item $T -> B: E_{K_{AT}}(K || R_A) || E_{K_{BT}}( K ||R_B) || N_A$
			\item $B -> A: E_{K_{AT}}(K || R_A) || E_{K}( R_A ||R_B) || N_A$
			\item $A -> B: E_{K}(R_B)$
		\end{enumerate}
	
		\textbf{Протокол с использованием хэш-функции}
		\begin{enumerate}
			\item $B -> A: R_B || text_1$
			\item $A -> B: text_3 || H_{K_{AB}}(R_A || R_B || ID_B || text_2)$
			\item $B -> A: text_5 || H_{K_{AB}}(R_B || R_A || ID_A || text_4)$
		\end{enumerate}
	}
	
	\frame{
		\frametitle{Протоколы "сильной" аутентификации}
		\framesubtitle{На основе асимметричных алгоритмов шифрования}
		
		\textbf{С использованием хэш-функции}
		\begin{enumerate}
			\item  $B -> A: h(R_B) || ID_B || E_{K_A}^{pub}(R_B || ID_B)$
			\item  $A -> B: R_B$
		\end{enumerate}	
	
		\bigskip
		
		
		\textbf{NSPK}
		\begin{enumerate}
			\item  $A -> B: E_{K_B^{pub}}(R_A || ID_A)$
			\item  $B -> A: E_{K_A^{pub}}(R_A || R_B)$
			\item  $A -> B: E_{K_B^{pub}}(R_B)$
		\end{enumerate}	
	}


	\frame{
		\frametitle{Протоколы "сильной" аутентификации}
		\framesubtitle{На основе асимметричных алгоритмов шифрования}
		
		\textbf{Атака параллельного сеанса}
		\textbf{NSPK}
		\begin{enumerate}
			\item[1]   $A -> I: E_{K_I^{pub}}(R_A || ID_A)$
			\item[1']  $I(A) -> B: E_{K_B^{pub}}(R_A || ID_A)$
			
			\item[2]   $B -> I(A): E_{K_A^{pub}}(R_A || R_B)$
			\item[2']  $I -> A: E_{K_A^{pub}}(R_A || R_B)$
			
			\item[3]   $A -> I: E_{K_I^{pub}}(R_B)$
			\item[3']  $I(A) -> B: E_{K_B^{pub}}(R_B)$
			
		\end{enumerate}	
	
		\textbf{Защита}
		\begin{itemize}
		
		\item
			\begin{enumerate}
				\item[2]  $B -> A: E_{K_A^{pub}}(R_A || R_B  || ID_B)$
			\end{enumerate}	
		\item
			\begin{enumerate}
				\item[3]  $A -> B: E_{K_{h(R_B)}}(ID_B)$
			\end{enumerate}	
			
		\end{itemize}
	}
	
	\frame{
		\frametitle{Протоколы "сильной" аутентификации}
		\framesubtitle{С использованием ЭЦП}
		
		\begin{enumerate}
			\item  $A -> B: cert_A || T_A / N_A || ID_B || text_2 || sign_A(T_A/N_A || ID_B || text1) $
			\item  $B -> A: cert_B || T_B / N_B || ID_A || text_4 || sign_B(T_B/N_B || ID_A || text3) $
		\end{enumerate}	
		
		\bigskip
		

		\begin{enumerate}
			\item  $B -> A: R_B || text_1$
			\item  $A -> B: cert_A || R_A || R_B || ID_B || text_3 || sign_A(R_A || R_B || ID_B || text_2)$
			\item  $B -> A: cert_B || R_B || R_A || ID_A || text_5 || sign_B(R_B || R_A || ID_A || text_4)$
		\end{enumerate}	
	}

	\frame{
		\frametitle{Протоколы "сильной" аутентификации}
		\framesubtitle{С использованием ЭЦП}
	
		
		
		\textbf{Атака}
		\begin{enumerate}
			\item  $I(B) -> A: R_B$
			\item  $A -> I(B): cert_A || R_A || R_B || ID_B || sign_A(R_A || R_B || ID_B)$
			
			\item[1']  $I(A) -> B: R_A$
			\item[2']  $B -> I(A): cert_B || R_B' || R_A || ID_A || sign_B(R_B' || R_A || ID_A )$
			
			
			\item  $I(B) -> A: cert_B || R_B' || R_A || ID_A || sign_B(R_B' || R_A || ID_A )$
		\end{enumerate}	
	}

	\frame {
		\frametitle{}
		
	}


	\frame {
	\begin{figure}
		%\includegraphics[width=0.8\linewidth]{index1.jpeg}
		
	\end{figure}
	}

\end{document}